% === CUMCM/MCM-ICM 数模论文完整 LaTeX 模板 ===
% 使用: xelatex paper.tex
\documentclass[12pt,a4paper]{article}
\usepackage[UTF8]{ctex}
\usepackage{amsmath,amssymb,amsfonts}
\usepackage{geometry}
\usepackage{graphicx}
\usepackage{booktabs}
\usepackage{algorithm}
\usepackage{algpseudocode}
\usepackage{hyperref}
\usepackage{cite}
\usepackage{array}
\usepackage{multirow}
\usepackage{xcolor}
\usepackage{float}
\usepackage{fancyhdr}
\usepackage{titlesec}
\usepackage[justification=centering]{caption}
\usepackage{setspace}
\usepackage{listings}
\usepackage{enumitem}
\usepackage{indentfirst}

\geometry{left=2cm,right=2cm,top=2cm,bottom=2cm}

% 页眉页脚
\pagestyle{fancy}\fancyhf{}\cfoot{\thepage}
\renewcommand{\headrulewidth}{0pt}
\fancypagestyle{plain}{\fancyhf{}\renewcommand{\headrulewidth}{0pt}\cfoot{\thepage}}

% 字体与行距
\setstretch{1.2}
\zihao{-4}

% 标题格式：一级居中，二三级左对齐
\titleformat{\section}
  {\centering\zihao{4}\heiti\bfseries}{\thesection}{1em}{}
\titleformat{\subsection}
  {\zihao{-4}\heiti\bfseries}{\thesubsection}{1em}{}
\titleformat{\subsubsection}
  {\zihao{-4}\heiti\bfseries}{\thesubsubsection}{1em}{}
\titleformat{name=\section,numberless}
  {\centering\zihao{4}\heiti\bfseries}{}{0em}{}

% 公式/图表编号
\numberwithin{equation}{section}
\numberwithin{figure}{section}
\numberwithin{table}{section}
\captionsetup{labelsep=quad,font=small,justification=centering}
\captionsetup[algorithm]{justification=centering}

% 代码格式
\lstset{basicstyle=\zihao{5}\ttfamily,numbers=left,numberstyle=\tiny,
  breaklines=true,frame=single,keepspaces=true,columns=flexible}

% 超链接
\hypersetup{colorlinks=true,linkcolor=black,citecolor=black,urlcolor=blue,bookmarks=true}

% 图形路径
\graphicspath{{data/figures/}{../figures/}{../../data/figures/}}

\title{}
\author{}\date{}

\begin{document}

% === 第一页：摘要页 ===
\thispagestyle{empty}

% 参赛编号
\begin{flushright}
\textbf{参赛编号:} \hspace{3cm}
\end{flushright}

\vspace{0.5cm}

% 论文题目（三号黑体，居中）
\begin{center}
\heiti\zihao{-2}
基于XXX模型的XXX问题研究
\end{center}

\vspace{0.8cm}

% 摘要正文
\begin{center}
\heiti\zihao{4}\textbf{摘要}
\end{center}

\vspace{0.3cm}

本文针对XXX问题，建立了XXX模型，求解得到XXX结果。

\textbf{针对问题一，}主要解决XXX问题。通过XXX分析，建立XXX模型，基于XXX方法/软件，求解得到XXX结果。

\textbf{针对问题二，}主要解决XXX问题。通过XXX分析，建立XXX模型，基于XXX方法/软件，求解得到XXX结果。

\textbf{针对问题三，}主要解决XXX问题。通过XXX分析，建立XXX模型，基于XXX方法/软件，求解得到XXX结果。

该模型具有XXX优点，可推广至XXX领域。

\vspace{0.5cm}

% 关键词（左对齐，3-5个，分号分隔）
\noindent\textbf{关键词：}XXX；XXX；XXX；XXX

\newpage

% === 目录（可选）===
% 2024年多数展示论文不设独立目录页。若需要目录，取消下面注释：
% \thispagestyle{empty}
% \tableofcontents
% \newpage

% === 正文开始 ===
\setcounter{page}{1}
\pagestyle{fancy}

% 1. 问题重述
\section{问题重述}
\subsection{问题背景}
（此处用自己的话描述赛题背景，不复制原文——说明情景、列出关键数据。）

\subsection{问题重述}
（此处列出题目要求解决的具体问题，包含题目给出的关键数字。⚠查重高风险区！）

% 2. 问题分析
\section{问题分析}
\subsection{问题一分析}
（问题一的数学本质是什么？为什么选择后续的建模方法？）

\subsection{问题二分析}
（问题二与问题一的依赖关系？方法选择理由？）

\subsection{问题三分析}
（以此类推，每子问题独立成小节。）

% 技术路线图（放在问题分析末尾）
% \begin{figure}[h]\centering
% \includegraphics[width=0.92\textwidth]{figures/技术路线图.png}
% \caption{总体技术路线}\label{fig:roadmap}
% \end{figure}

% 3. 模型假设
\section{模型假设}
\textbf{假设1：}（具体假设内容，含关键处理细节和量化数据。）\\
\textbf{假设2：}（每条假设解释为什么做此假设，与后续模型有何关联。）\\
\textbf{假设3：}（禁止空洞假设——不写"仅考虑理想条件"一类套话。）

% 4. 符号说明
\section{符号说明}
（此处用三线表列出主要全局变量——三列：符号、含义、单位。临时变量不列。）

% 5. 问题一模型的建立与求解
\section{问题一模型的建立与求解}
\subsection{数据预处理}
（数据来源、清洗逻辑、关键数字。用文字说明，不是只列公式。）

\subsection{模型建立}
（建模动机→逐步推导→公式前后有文字解释→约束条件含义说明。）

\subsection{模型求解}
（算法步骤、参数设置。复杂算法配伪代码或流程图。）

\subsection{结果分析}
（表格+文字解读+图表引用。数值必须与代码输出一致。）

% 6. 问题二模型的建立与求解（同上结构）
% 7. 问题三模型的建立与求解（同上结构）

% 灵敏度分析（独立成章）
\section{灵敏度分析}
\subsection{模型对XXX参数的灵敏性}
（参数在[范围]内以步长变化 → 目标值变化 → 敏感/不敏感？结论。）

% 模型检验（三种形式：稳定性敏感性分析 / 统计检验误差分析 / 新旧模型对比）
\section{模型检验}
（Q1预测精度验证(WAPE/MAPE/RMSE+分层)、Q2约束可行性验证(每类约束违反数)、
Q3方案验证(重优化违反数/成本口径)、基准对比、综合PASS声明。）

% 模型评价与推广（三段式：优点→缺点→改进推广）
\section{模型评价与推广}
\subsection{模型的优点}
（3-4条，每条对应具体方法贡献——禁止"模型合理""结果可靠"。）
\subsection{模型的缺点}
（2-3条诚实局限——按重要性排序。）
\subsection{模型的改进与推广}
（改进方向具体可操作，推广说明适用范围。）

% === 参考文献（GB/T 7714，完整8种文献类型） ===
\section{参考文献}
\begin{thebibliography}{99}
\bibitem{ref1} 作者. 题名[J]. 期刊名, 年, 卷(期): 页码.
\bibitem{ref2} 作者. 书名[M]. 出版地: 出版社, 出版年.
\bibitem{ref3} 作者. 题名[D]. 出版地: 出版者, 出版年.
\bibitem{ref4} 作者. 题名[C]//论文集名称. 会议地点, 会议日期.
\bibitem{ref5} 专利申请者. 专利题名: 专利国别, 专利号[P]. 公告日期.
\bibitem{ref6} 作者. 题名[EB/OL]. (发布日期)[引用日期]. 获取路径.
\bibitem{ref7} 作者. 题名[R]. 出版地: 出版者, 出版年.
\bibitem{ref8} 标准号. 题名[S]. 出版地: 出版者, 出版年.
\end{thebibliography}

% === 附录 ===
% 附录内容（不限六件套，还可包括：详细证明、外部数据、大流程图、繁杂图表）
\appendix
\renewcommand{\thesection}{附录\Alph{section}}

% 附录A：支撑文件列表
\section*{附录A\quad 支撑文件列表}
\addcontentsline{toc}{section}{附录A\quad 支撑文件列表}
支撑文件目录结构：
\begin{verbatim}
├── Q1/solve_q1.py, plot_q1.py
├── Q2/milp_model.py, verify_q2.py
└── result1.xlsx, result2.xlsx
\end{verbatim}

% 附录B~D：各问题求解结果表
\section*{附录B\quad 问题一求解结果}
\addcontentsline{toc}{section}{附录B\quad 问题一求解结果}
% 完整数表（正文引用："详见附录B"）

\section*{附录C\quad 问题二求解结果}
\addcontentsline{toc}{section}{附录C\quad 问题二求解结果}

\section*{附录D\quad 问题三求解结果}
\addcontentsline{toc}{section}{附录D\quad 问题三求解结果}

% 附录E：关键公式推导
\section*{附录E\quad 关键公式推导}
\addcontentsline{toc}{section}{附录E\quad 关键公式推导}
% 正文省略的长推导

% 附录F：核心程序代码
\section*{附录F\quad 核心程序代码}
\addcontentsline{toc}{section}{附录F\quad 核心程序代码}
\lstset{basicstyle=\zihao{5}\ttfamily, frame=single, breaklines=true, numbers=left}
\begin{minipage}{\textwidth}
\lstinputlisting[language=Python]{code/solve.py}
\end{minipage}

\end{document}
